\documentclass[12pt]{article}
\usepackage[utf8]{inputenc}
\usepackage{amsmath, amssymb, physics, graphicx, tikz}
\usepackage[a4paper, margin=1in]{geometry}
\usepackage{fancyhdr}
\usepackage{hyperref}
\usepackage{tikz}
\usepackage{xcolor}
\definecolor{sectionrectanglecolor}{rgb}{0.25,0.5,0.75}
\definecolor{sectiontitlecolor}{rgb}{0.2,0.4,0.65}
\definecolor{subsectioncolor}{rgb}{0,0,0}
\definecolor{lightgray}{gray}{0.9}
\definecolor{darkgray}{gray}{0.1}
\definecolor{lightred}{rgb}{0.9,0,0.1}
\usepackage{wrapfig}
\usepackage{tcolorbox}

\newcommand{\be}{\begin{eqnarray}}
\newcommand{\non}{\nonumber \\ }
\newcommand{\ee}{\end{eqnarray}}
\newif\ifshowboardanswers
\showboardanswersfalse
% Set this to \showboardquestiontrue to include the end-of-lecture board question.
\newif\ifshowboardquestion
\showboardquestionfalse
\newcommand{\boardanswer}[1]{\ifshowboardanswers\\[0.35em]\textcolor{blue}{\textbf{Answer.} #1}\fi}

\fancyhead[L]{Special Relativity}
\fancyhead[R]{Lecture 5}

\title{\textbf{Lecture 5: The Doppler effect and Rapidity}}
\author{Based on Morin's \textit{Introduction to Classical Mechanics}, Ch. 11.8--11.9}
\date{}

\begin{document}
\maketitle

\section*{Outline}
\begin{enumerate}
    \item \href{sec:recap}{Recap lecture 4}
    \item \href{sec:DE}{The Doppler effect}
    \item \href{sec:Rap}{Rapidity}
\end{enumerate}

\section{Recap Lecture 4}\label{sec:recap}
We showed that we can define a frame-independent quantity, the spacetime interval (or interval, according to Morin):
\begin{tcolorbox}[colback=blue!5!white]
\vspace{-1.3em}
\be
\Delta s^2  = c^2 \Delta t^2 - \Delta d^2 \label{eq:SI}.
\ee
\end{tcolorbox}
We showed that $c\Delta t \geq \Delta s$. The formal definition of the timelike spacetime interval states that it is $c$ times the time between two events as measured by an inertial clock present at both events. Because the inertial clock follows a unique worldline, the spacetime interval is unique for a given pair of events (hence it is frame-independent). A geometrical analogy is the distance in Euclidean space.

Besides the spacetime interval, we also derived an expression for proper time:
\begin{tcolorbox}[colback=blue!5!white]
\vspace{-1.3em}
\be
\Delta \tau = \int_{t_A}^{t_B} \sqrt{1-\frac{|\vec{v}|^2}{c^2}}\,dt.
\ee
\end{tcolorbox}
The formal definition for proper time is the time between two events as measured by a single clock present at both events. Since this clock does not have to be inertial, $c\Delta\tau$ is not equal to the spacetime interval, unless that single clock follows an inertial path (i.e., a straight line in a spacetime diagram). The geometrical equivalent is the path length: the time measured on the clock depends on the precise worldline the clock follows between the two events. While this is a geometrical analogy, the hierarchy is the opposite $c\Delta t\geq \Delta s \geq c\Delta \tau$.

Next, we introduced spacetime diagrams. They allow us to improve our intuition for simple 2D problems. For SR diagrams, it does help to work in SR units (although it is possible to simply label the time axis as $ct$ as Morin does). In that case, we can treat time and space on equal footing and use units that are identical on both axes. Two observer diagrams are useful when considering scenarios where we have to compare two frames. We derived how to draw the axes of the frame $S'$ moving with speed $u$ in the positive $x$ direction, with $\beta_u\equiv u/c$, and how to put ticks on these axes. We showed how we can use the two observer diagram to understand length contraction.

In this lecture we continue this convention: $u$ is the speed of the boost from $S$ to $S'$, $\beta_u\equiv u/c$, and $\gamma_u=(1-\beta_u^2)^{-1/2}$. We reserve $v$ (and $v'$) for the speed of a physical object measured in a specified frame.

\subsection{The pole and barn paradox}
To show how two observer diagrams can help understand perhaps contradicting statements, we will consider the pole and barn paradox. In this hypothetical scenario, we consider two observers. One observer runs with a pole of proper length $L'_p = 10$ light nanoseconds (lns, SR units) past a barn of proper length $L_{\rm barn} = 8$ lns\footnote{Recall, proper length is the length of an object measured at rest.}. The runner's frame $S'$ moves at $u=3c/5$ relative to the barn frame $S$ (pretty fast, yes, it is a paradox), so $\gamma_u=5/4$. The barn observer finds the pole contracted: $L_{p}=\gamma_u^{-1}L'_p=8$ lns. The pole fits {\it in} the barn. The runner, on the other hand, sees the barn contracted: $L'_{\rm barn}=\gamma_u^{-1}L_{\rm barn}=6.4$ lns. The runner concludes that the pole does not fit in the barn! Who is right?

The solution to this paradox requires defining length in terms of events! In that case, there is no such general thing as the pole `fitting’ the barn because that relies on events being simultaneous (which depends on coordinate time and thus on the reference frame).

The solution can be understood if we draw a two observer diagram. We show the steps how to draw such a diagram in fig.~\ref{fig:PBP}. The main point is that both are right from the perspective of their frames, since what is simultaneous in one frame is not simultaneous in another. Since an object's length in a given inertial frame is defined to be the distance between any two simultaneous events that occur at its ends, they will not agree, but that is fine and fully consistent with SR.

There is a large number of problems that are better understood if you draw a (two observer) spacetime diagram. Morin has many questions, and I encourage you to take a look. For example, you can also derive the rear clock ahead effect using a spacetime diagram. You can find the solution to this problem in one of the exams.

An even more challenging paradox is the skateboard and pothole example. In that example, there is a skateboard (1lns) moving to the right, and a pothole moving up (1lns). Does the skateboard fall through the pothole? If you send me a solution, you will receive a reward.

\begin{figure}
    \centering
    \includegraphics[width=1\linewidth]{Figures/PBP.pdf}
    \caption{{\bf Top left:} two observer diagram and the worldregion of the barn (8ns) in the frame at rest. {\bf Top Right:} adding $x'$ and $t'$ coordinate axes for the $S'$ frame moving with $3/5 c$. {\bf Bottom left:} the pole fits exactly in the barn according to the observer at rest, which is clear from the intersection of the two worldregions (barn and pole) as measured along the spatial coordinate axis $x$. {\bf Bottom right:} the runner disagrees. In the runner's frame, simultaneous events are parallel to the $x'$ axis, and from $S'$ the two events that are simultaneous show that the pole does not fit in the barn (the pole's front exits the barn before the rear enters the barn, which defines the length of the barn compared to the pole). }
    \label{fig:PBP}
\end{figure}

\section{The Doppler effect}\label{sec:DE}
\subsection{Longitudinal effect}
In previous lectures, we have discussed the difference between seeing and observing. The latter refers to when something has actually happened, as can be read off from a lattice with a set of synchronized clocks (the value of $t$ on the coordinate axis). Seeing refers to when the light that comes off an event actually hits your eye. The Doppler effect cares about seeing.

Consider a source that flashes at frequency $\nu'$ in its own rest frame $S'$. Relative to the observer frame $S$, the source frame moves away with constant speed $u$, so $\beta_u=u/c$.
The relativistic Doppler effect has two contributions:
\begin{itemize}
\item Time dilation: There is more time between flashes in your frame ($S$), which means that the frequency $\nu$ in your frame should go down.
\item Classical (geometric) Doppler: similar to sound, this depends on the component of the source velocity along the line of sight. If a source moves towards you, its distance decreases and successive wavefronts arrive closer together; if it moves away, its distance increases. This increases the observed frequency for approach and reduces it for recession. A passing siren consequently changes from high to low pitch, but at the instant of closest approach its motion is transverse to the line of sight and the instantaneous classical shift is zero.
\end{itemize}
To obtain an actual expression for the Doppler effect, we first note that the time interval between the emissions as measured in frame $S'$ can be expressed as $\Delta t' = 1/\nu'$, i.e., one over the frequency of the flashes as measured in frame $S'$. We can apply time dilation to find that the time interval between emissions in frame $S$ is given by $\Delta t = \gamma_u \Delta t'$. Now this is not the time between the flashes as seen by the observer, since the source is moving and the distance to it changes between emissions. Specifically, the source has travelled $u\Delta t=u\gamma_u\Delta t'$ in between emissions. So the time $\Delta T$ between flashes as seen by the observer in $S$ is given by the difference in travel distance between the two emissions, divided by the speed of a flash ($c$), i.e.
\be
\Delta T = \frac{1}{c}(c+u)\gamma_u \Delta t' = \frac{1+\beta_u}{\sqrt{1-\beta_u^2}} \Delta t' = \sqrt{\frac{1+\beta_u}{1-\beta_u}} \nu'^{-1}, \label{eq:Doppler}
\ee
The frequency of the observed source is therefore
\begin{tcolorbox}[colback=blue!5!white]
\vspace{-1.3em}
\be
\nu = \sqrt{\frac{1-\beta_u}{1+\beta_u}} \nu'.
\ee
\end{tcolorbox}
\noindent
If a source is moving away from you, $\beta_u>0$ and $\nu<\nu'$. The light of the source is redshifted. If the source is moving towards you, replace $\beta_u$ by $-\beta_u$, and the light is blueshifted. We show a spacetime diagram in Fig.~\ref{fig:Doppler} sketching the scenario for a source moving away. The figure is using SR units, so $c = 1$ (and hence the total time between flashes $d t_R$ misses factors of $c$). To obtain the observed frequency between flashes, one needs to relate $dt_R=\gamma_u\Delta t_R'$.


\begin{figure}[h]
    \centering
    \includegraphics[width=\linewidth]{Figures/Doppler.pdf}
    \caption{{\bf Left:} Relativistic Doppler effect for a source emitting flashes at equal time intervals in frame $S'$. The flashes are seen by an observer in frame $S$ at intervals that are either blueshifted (when the source is moving towards the observer) or redshifted (when the source is moving away, as shown in this diagram). {\bf Right:} Doppler effects when the source is not exactly moving towards or away from the observer. With positive $\beta_u$ denoting recession, the longitudinal factor becomes $1+\beta_u\cos\Theta$.}
    \label{fig:Doppler}
\end{figure}


\subsection{Transverse effect}
The expression we derived in the previous section holds for sources that are either moving directly towards you or directly away from you, the so-called longitudinal Doppler effect. What happens when a source is not moving exactly along the line of sight, but at an angle $\Theta$ (see the right side of Fig.~\ref{fig:Doppler})? The changing-distance contribution then involves the projected velocity, so $1+\beta_u\rightarrow 1+\beta_u\cos\Theta$. Inserting that into Eq.~\eqref{eq:Doppler}, we find
\be
\nu  = \gamma_u^{-1} (1+\beta_u \cos \Theta)^{-1} \nu'.
\ee
We can consider the extreme case when the source is moving transverse to the observer. In that case $\Theta = \pi/2$ and we find
\be
\nu = \nu' \sqrt{(1-\beta_u^2)}.
\ee
Note that this is a second-order effect. At the instant when a classical sound source moves transversely to the line of sight, its radial velocity is zero, so there is no instantaneous classical Doppler shift. For light, however, the source clock still runs slowly, which produces the transverse redshift. The above scenario is equivalent to `Case 2' in Morin.
%[Have to think about a way to explain case 1; perhaps it is not possible with this simple example. ]

\section{Rapidity}\label{sec:Rap}
Recall that we could write the Lorentz transformations in matrix notation:
\be
\begin{pmatrix}
c\Delta t' \\
\Delta x'
\end{pmatrix} =
\begin{pmatrix}
\gamma_u & -\beta_u \gamma_u \\
-\beta_u \gamma_u  & \gamma_u
\end{pmatrix}\begin{pmatrix}
c\Delta t \\
\Delta x
\end{pmatrix}.
\ee
Now, let us define the rapidity $\phi_u$ of the boost as
\be
\tanh \phi_u \equiv \beta_u,
\ee
There are some very nice properties of this function, which simplify SR transformations.

For example, consider Einstein’s velocity transformations:
\be
\beta_v\equiv\frac{v_x}{c}=\frac{\beta_{v'}+\beta_u}{1+\beta_u\beta_{v'}}=\frac{\tanh\phi_{v'}+\tanh\phi_u}{1+\tanh\phi_u\tanh\phi_{v'}}=\tanh(\phi_{v'}+\phi_u), \label{eq:RV}
\ee
where the equality between the last two expressions is the consequence of the properties of the hyperbolic trig function (appendix A of Morin).

Similarly:
\be
\gamma_u \equiv \frac{1}{\sqrt{1-\beta_u^2}} = \frac{1}{\sqrt{1-\tanh^2 \phi_u}} = \cosh \phi_u,
\ee
and
\be
\gamma_u \beta_u = \sinh \phi_u.
\ee
This allows us to write the Lorentz transformation as
\be
\begin{pmatrix}
c\Delta t' \\
\Delta x'
\end{pmatrix} =
\begin{pmatrix}
\cosh \phi_u & -\sinh \phi_u \\
-\sinh \phi_u  & \cosh \phi_u
\end{pmatrix}\begin{pmatrix}
c\Delta t \\
\Delta x
\end{pmatrix},
\ee
which is remarkably similar to rotations around an angle $\alpha$ in Euclidean space:
\be
\begin{pmatrix}
x' \\
y'
\end{pmatrix} =
\begin{pmatrix}
\cos \alpha &  \sin \alpha \\
-\sin \alpha & \cos \alpha
\end{pmatrix}
\begin{pmatrix}
x \\
y
\end{pmatrix}.
\ee
Instead of trig functions for rotations, we have hyperbolic trig functions.

How can this be useful? Suppose, for example, that you apply Lorentz transformations twice (see Q1 R3, 2023), once with boost speed $u_1$ (associated with rapidity $\phi_{u_1}$) and another with boost speed $u_2$ (associated with rapidity $\phi_{u_2}$). Doing this without rapidity is a bit of a mess, but it should lead to another Lorentz transformation, now with the boost speed given by the Einstein velocity transformation. With rapidity, this follows immediately by applying the Lorentz-transformation matrix twice and using hyperbolic trigonometric identities: $\phi_{u,\rm final}=\phi_{u_1}+\phi_{u_2}$ and $u_{\rm final}/c=\tanh(\phi_{u_1}+\phi_{u_2})$, identical to Eq.~\eqref{eq:RV}.

\subsection{Interpretation}
See Morin. When going from Eq.~(11.62) to Eq.~(11.63), you have to do two things. First, we should realize that locally, any function can be written as the function plus an infinitesimal displacement, which is determined by the first derivative multiplied by the displacement:
\be
x(t + d t) = x(t) + \frac{\partial x}{\partial t} d t + \mathcal{O}(d t^2).
\ee
Here $\partial x /\partial t$ is the partial derivative, the derivative of a function keeping all other parameters fixed. But since $v$ does not depend on anything but time, it can be replaced with the ordinary derivative below. To expand to first order in $dt$ requires expanding the RHS of Eq.~(11.62). We write:
\be
\frac{v(t) + a dt}{1+ v(t) a dt/c^2} &=& (v(t) + a dt)(1+v(t) a dt/c^2)^{-1}\nonumber \\
&=& (v(t) + a dt)(1 - v(t) a dt/c^2 + \mathcal{O}(dt^2)).
\ee
We thus find
\be
v(t) + \frac{dv}{dt} dt = v(t) + a dt( 1 - v(t)^2/c^2),
\ee
to first order in $dt$. This is equivalent to Eq.~(11.63).

\section{Relativity without $c$}
Please read this section in Morin. It is a very nice derivation and quite powerful. It will not be something we will put in an exam (perhaps only a conceptual question).


\ifshowboardquestion
\section*{Board Question}
\begin{tcolorbox}[colback=blue!5!white]
A source moves directly away from Earth with $\beta_u=0.60$ and emits light with wavelength $\lambda_{\rm em}=500\,{\rm nm}$ in its own rest frame.
\begin{enumerate}
    \item Compute the observed wavelength on Earth using the relativistic Doppler formula.
    \boardanswer{For recession, $\lambda_{\rm obs}=\lambda_{\rm em}\sqrt{(1+\beta_u)/(1-\beta_u)}=500\,{\rm nm}\sqrt{1.6/0.4}=1000\,{\rm nm}$.}
    \item Express the same boost using a rapidity $\phi_u$.
    \boardanswer{$\tanh\phi_u=\beta_u=0.60$, so $\phi_u=\tanh^{-1}(0.60)=\frac12\ln(4)=\ln 2$.}
    \item A second source moves away from the first with $\beta_{u_2}=0.60$. Use rapidities to find the speed of the second source relative to Earth.
    \boardanswer{Rapidities add: $\phi_{u,\rm tot}=2\ln 2=\ln 4$. Thus $\beta_{u,\rm tot}=\tanh(\ln 4)=15/17\simeq0.88$.}
    \item Explain why ordinary addition of the two speeds would give an impossible answer.
    \boardanswer{Ordinary addition gives $1.2c$, which exceeds the speed of light. In SR, rapidities add linearly while velocities combine nonlinearly.}
\end{enumerate}
\end{tcolorbox}
\fi

\end{document}
